\newpage
\section*{Agradecimientos}
\vspace{1cm}

Escribir este apartado supone algo más que unos simples agradecimientos; representa el cierre de una etapa de esfuerzo, aprendizaje y superación personal. La culminación de este TFG y, con él, de mi etapa universitaria, algo que al inicio parecía que no tener fin. Sin embargo, todo acaba llegando, y en el camino siempre aparecen personas que hacen que los retos sean más llevaderos. Personas sin las cuales no habría sido posible alcanzar esta meta y que, por ello, merecen un lugar especial en estas páginas.

En primer lugar, quiero expresar mi más sincero agradecimiento a mis tutores, David Guijo Rubio y Victor Manuel Vargas Yun. Habéis sido los más rápidos revisando mis entregas y encontrando fallos en aquello que yo creía perfecto. Gracias por estar al pie del cañon hasta el final y, sobre todo, por todos los consejos y correcciones que me habéis ido dando durante el proceso. Sin vuestra ayuda y dedicación habría sido imposible llegar hasta este punto.

A mi familia, especialmente a mis padres, Juan Pedro y Trini; a mi hermana, Graciela; y a mis abuelos, Estanislao y Juana. Gracias por siempre estar ahí, gracias por todo lo que me habéis enseñado desde pequeño y la educación que me habéis dado. Gracias por hacerme la persona que soy hoy en día y por enseñarme que el esfuerzo, la constancia y el trabajo acaban dando sus frutos. Siempre voy a estar innmensamente agradecido con vosotros y este documento es la prueba del buen trabajo que habéis hecho.

A Elena, porque eres quien mejor sabe todo lo que he pasado para llegar hasta aquí y porque sé que este logro te hace tanta ilusión como a mí. Gracias por acompañarme todo este tiempo, por escucharme, apoyarme y animarme siempre que lo he necesitado. Gracias por estar a mi lado y por celebrar conmigo cada pequeño paso que me ha traído hasta aqui, por eso sé que este trabajo también lleva una pequeña parte de ti.

A mi Carlos y Matas, por todos estos años compartiendo paredes, compartiendo vida, anécdotas y locuras. Habéis hecho que me sienta como en casa desde el primer día y habéis sido los mejores en enseñarme que no todo en la vida es organización, estudiar y perfección, que hay algo más importante que todo eso, la diversión. Gracias de corazón por estar siempre ahí, nunca voy a olvidar la cantidad de momentos que hemos vivido juntos y los que nos quedan por vivir.  

Al resto de mis cordobeses, Manuel, Juan, Ismael, Ruso, y a mi malagueño, Pepe, sin duda soy el más feliz del mundo con la familia que hemos formado durante estos años. Gracias por hacer esta piña tan increible, por amenizar las clases y por todos los momentos compartidos tanto dentro como fuera de la universidad. Espero que esto no termine aquí y que sigamos manteniendo la amistad que hemos construido. Os deseo lo mejor en lo que os queda por delante y, por supuesto, ¡nos vemos en Croacia!

A todos, de corazón, gracias.
